\section{Introduction}
\label{sec:intro}

A rectenna converts radiation to direct current by feeding an antenna into a
rectifier. Above a few hundred gigahertz the rectifier is the difficulty: the
$RC$ time of a junction diode limits Schottky rectifiers to the low terahertz
range, and the search for a junction-free, low-capacitance nonlinearity that
works to $1\THz$ and beyond is long standing. Graphene has attracted two
proposals of this kind, both of which rest on the shape of a channel rather
than on a junction.

The first is the geometric diode~\cite{moddel2012ultrahigh,zhu2013graphene,dragoman2013geometrically}: a channel
that narrows asymmetrically, typically a funnel, whose dimensions are below the
mean free path so that carriers are steered ballistically by the walls. With
planar geometry and negligible capacitance it was proposed explicitly for
terahertz and infrared rectennas~\cite{zhu2013graphene}. Encapsulated devices
reach reported asymmetry ratios of $1.6$ to
$2$~\cite{nguyen2021geometric,wang2024graphene}, and related four-terminal
ballistic rectifiers have been used for detection and imaging up to
$0.7\THz$~\cite{auton2016graphene,auton2017terahertz}. The second is the
electronic Tesla valve. Tesla's valvular conduit~\cite{tesla1920valvular}
rectifies fluid flow by inertia~\cite{nguyen2021early}, and in the temperature
window where electron-electron collisions make the graphene electron liquid
viscous~\cite{bandurin2016negative,krishnakumar2017superballistic} a Tesla-shaped
channel has been reported to show a five per cent asymmetry of its
four-terminal resistance~\cite{geurs2020rectification}. A shape that rectifies an electron
fluid with no junction is an obvious rectenna candidate.

What the shape can do is constrained from the outset. For an elastic
two-terminal conductor the unitarity of the scattering matrix implies equal
transmission in the two directions at every energy, so the linear response
is symmetric whatever the shape. Rectification is a second-order effect,
and a consistent theory of it requires the response of the potential
landscape inside the device to the
bias~\cite{christen1996gauge,sanchez2004magnetic,song1999formalism,lofgren2004symmetry}.
The useful question is therefore not whether a shape rectifies but which
nonlinearity it feeds, how strong that nonlinearity is at the power an antenna
delivers, and how fast it is. A geometric diode in the ballistic regime feeds
the bias-induced landscape; a Tesla valve in the viscous regime feeds the
convective term of the Navier-Stokes equation; and at terahertz frequencies a
gated graphene channel is a plasmonic cavity whose convective and
density nonlinearities rectify by the mechanism of Dyakonov and
Shur~\cite{dyakonov1993shallow,dyakonov1996detection}, which is the basis of
most graphene terahertz detectors to
date~\cite{vicarelli2012graphene,bandurin2018resonant}. These three
mechanisms have been studied separately, with different geometries and
different figures of merit. Whether a geometric shape, as opposed to the
terminations of the channel, carries the terahertz response has not been
tested with one device in all of them.

This paper does that. We take one set of geometries
(Fig.~\ref{fig:devices}): the funnel of the best-documented geometric
diode~\cite{nguyen2021geometric}, the classical Tesla valve at the same
minimum feature size, a straight control, and versions scaled to a $1\um$
channel that a $75\ \Omega$ antenna can drive. We compute their rectification
(i) in a semiclassical ballistic model of the kind used in the geometric-diode
literature~\cite{zhu2013graphene,stearns2021simulation,truccolo2024comprehensive}, (ii) in
steady viscous electron hydrodynamics, (iii) in the gated compressible
(plasmonic) electron fluid from $10\GHz$ to $3\THz$, at small signal and at
finite amplitude up to the gate overdrive, and (iv) in the ballistic model
driven at terahertz frequencies. Each solver is validated against analytic
limits and benchmarks before it is used; the acceptance tests are listed in
the Supplemental Material.

The answer is the same in every regime, and it is negative for the shape.
The Tesla valve does not rectify ballistically, rectifies viscously only at
Reynolds numbers the Fermi liquid cannot sustain, and contributes
$10^{-5}$ to $10^{-3}\AW$ at terahertz frequencies. The funnel rectifies in
all regimes, but its terahertz response of $0.1$ to $0.3\AW$ is the
Dyakonov-Shur asymmetry of its terminations; the shape adds a factor of
order one near a taper resonance and nothing elsewhere. The conversion
efficiency of the plasmonic element is proportional to the incident power up
to a saturation power set by the gate, and its peak is below one per cent at
every operating point for which the fluid description holds. The ballistic
shape asymmetry, driven at $1\THz$, falls to the same responsivity. As a
detector the element is respectable; as a harvester it does not exist.

Section~\ref{sec:models} describes the devices and the three models,
Sec.~\ref{sec:dc} the dc results, Sec.~\ref{sec:plasmonic} the plasmonic
element from small signal to saturation, Sec.~\ref{sec:ballistic} the
ballistic element at terahertz frequencies, and Sec.~\ref{sec:discussion}
compares the regimes, places the numbers against measured detectors and
states the limitations.

\begin{figure*}
\includegraphics[width=\textwidth]{\figpath/fig1_devices.pdf}
\caption{Devices and frequency scales. (a) The funnel geometric diode,
device D of Ref.~\cite{nguyen2021geometric}: $2\um$ leads, a 30 degree taper
to a $175\nm$ neck and an abrupt shoulder; forward is the direction into the
taper. (b) The classical four-loop Tesla valve at the same $175\nm$ channel
width. (c) The gated funnel scaled to a $1\um$ neck with $4\um$ contacts, with
the terminations used for the plasmonic response: the antenna, a Thevenin
source of resistance $\Rant$, drives the source contact and the drain is tied
to the gate. (d) Characteristic frequencies of the mechanisms studied, for the
parameters of Table~\ref{tab:params}; the shaded band is the range of
interest for a terahertz rectenna.}
\label{fig:devices}
\end{figure*}
